% =====================================================================
%  Preguntas abiertas al comité
%  Anexos
%
%  ENCARGO: Las decisiones que dependen de los asesores y su estado.
%
%  Reglas del documento:
%   - Una idea = una subseccion = una figura.
%   - Toda figura declara su procedencia con.
%   - Toda cifra sale del canon de agosto; ver GUIA.md.
%   - M1 y M3 se reportan siempre en paralelo.
%   - Ningun superlativo sin releer el CSV de la frontera correcta.
% =====================================================================
\section{Preguntas abiertas al comité}
\label{anx:preg}
\justifying

Este anexo reúne las decisiones que no corresponden a este trabajo y que se
dejan planteadas con la evidencia reunida, para que quien decida lo haga
sobre cifras y no sobre impresiones. Cada pregunta se enuncia, se acompaña
de lo medido y termina con lo que cambiaría según la respuesta.

\subsection{Las fuentes del modelo base y cuál manda}
\label{sub:preg-competencia}

Esta entrada nació como pregunta y dejó de serlo al aparecer la versión
arbitrada del modelo. Se conserva porque el desacuerdo entre fuentes sigue
existiendo y conviene que quede constancia de cuál se siguió.

El bienestar del comprador incluye un término que penaliza competir con los
demás, y un segundo término que recoge lo que paga por la energía. El guion
en Python que acompaña al modelo escribe el primero multiplicando por la
energía del comprador propio, mientras que la versión arbitrada, publicada
en \emph{IEEE Latin America Transactions} en agosto de 2025, lo escribe
multiplicando por la energía del comprador ajeno, igual que la formulación
que aparece comentada en el código fuente en Matlab. Para el segundo, el
documento extenso del modelo introduce un factor del precio de compra a la
red que ni la versión arbitrada ni el guion contemplan.

Las dos expresiones del primer término coinciden exactamente cuando todos
los compradores adquieren la misma cantidad, que es el caso de las
verificaciones sintéticas con que se validó la traducción en su momento, y
se separan en cuanto las cantidades difieren. La magnitud de la separación
se midió sobre cuatro compradores con cantidades distintas: la diferencia
mayor asciende a \num{1755,95} unidades sobre valores del orden de
\num{2500}.

\textbf{Este trabajo sigue en los dos casos a la versión arbitrada}, que es
la especificación publicada. La elección está además aislada: se comprobó
que ninguna de las funciones que resuelven el mercado consulta esas
expresiones, de modo que no mueve ni un flujo ni un precio. De los \num{672}
valores que componen la huella de verificación de una jornada completa,
cambiar de forma altera \num{25}, y los veinticinco son bienestar del
comprador.

Queda por confirmar, y es lo único que sigue abierto en este punto, si el
guion en Python recoge una versión anterior del modelo o una corrección
posterior que no llegó al artículo.

\subsection{Las dos preguntas regulatorias}
\label{sub:preg-regulatorias}

Quedan planteadas y sin resolver desde el repositorio.

La primera es si la exención de la contribución de solidaridad está
efectivamente radicada para las cinco instituciones, o si es solo aplicable
en derecho. El trabajo la aplica a las cinco, por tratarse de usuarios no
regulados y de servicios de salud y educación, pero eso es una lectura de la
norma y no la verificación de un trámite.

La segunda es si la energía intercambiada dentro de la comunidad paga o no
cargos por uso de redes. La respuesta gobierna el ancho de la banda de
precios y, con él, el valor entero del mercado: sin exención la banda se
vacía en la gran mayoría de las horas.

Para esta segunda conviene llevar un antecedente del propio grupo. El
trabajo sobre mercados entre pares en escenarios multimicrorred, publicado
en \emph{TecnoLógicas} en 2024, modela explícitamente una penalización, o
costo de transmisión, cuando un recurso envía energía a una microrred
vecina. Es decir, el grupo ya ha modelado que el transporte entre zonas se
cobra, lo que sitúa la pregunta en un terreno conocido en lugar de abrirla
desde cero.

\subsection{Si el techo del precio es uno solo o uno por comprador}
\label{sub:preg-techo}

El límite superior del precio del mercado es el costo unitario que cada
institución le paga a su comercializador, porque por encima de él le
conviene comprarle a la red. La comunidad tiene dos comercializadores, de
modo que ese límite no es el mismo para las cinco: cuatro instituciones lo
tienen en 731,1 (COP/kWh) y la quinta en 777,2.

La consulta al asesor externo en materia regulatoria, de junio de 2026,
parece resolver el asunto y no lo resuelve. Su respuesta afirma tres veces
que el costo unitario base es el mismo para todos, pero contesta por clase
de usuario, es decir comercial frente a oficial dentro de un mismo
comercializador, y esa distinción ya se retiró al reconocer a las cinco como
usuarios no regulados. Lo que hoy separa los dos techos es el
comercializador, y sobre ese eje no se preguntó.

Se midieron los dos regímenes sobre sesenta horas activas de cada frontera,
evaluando siempre el ahorro contra el costo unitario real de cada
institución. La energía transada es idéntica en ambos, con diferencias del
orden de una diezbillonésima de kilovatio hora, lo que confirma que el
límite superior no interviene en la cantidad. Lo que cambia es el
emparejamiento y el reparto.

Con el techo de cada comprador, el excedente de la frontera principal
asciende a 46.596,6 (COP) y el 82 \% del ahorro del lado comprador se
concentra en la institución del otro comercializador, porque la energía
escasa va hacia quien tiene la alternativa externa más cara. Con un techo
único e igual al más bajo, el excedente se sitúa en 45.016,1 (COP), un
3,4 \% menos, y el ahorro se reparte de forma más pareja. La disyuntiva, por
tanto, es entre eficiencia y equidad, y no entre correcto e incorrecto.

Un tercer régimen queda descartado sin necesidad de decisión: un techo único
e igual al más alto obliga a las cuatro instituciones del techo bajo a pagar
756,1 (COP/kWh) por energía que la red les vendería a 731,1, con lo que su
ahorro se vuelve negativo en las ochenta y ocho horas medidas. Ningún
régimen que obligue a un participante a pagar por encima de su alternativa
externa es admisible.

La implementación admite las dos respuestas. Cuando los techos coinciden, la
formulación por comprador se reduce exactamente a la de techo único, de
manera que la decisión del comité no obliga a reescribir nada.

\subsection{Cómo se liquida un mercado horario con un crédito mensual}
\label{sub:preg-liquidacion}

Esta entrada nació como pregunta y la resolvió el propio texto regulatorio. El
Anexo~4 de la Resolución CREG 101 072 de 2025, que aplica el artículo 26 de la
Resolución CREG 174 durante el periodo de transición, liquida el crédito con
las cantidades del mes y valora a la bolsa de cada hora lo que lo supera, a
partir de una hora de corte que define como aquella en que los excedentes
horarios acumulados desde el inicio del mes igualan o sobrepasan la
importación total del mes. El cupo se fija con el total mensual, la hora del
corte con el acumulado horario, y el dinero se liquida al cierre del período.

Un mercado entre pares encaja en esa misma estructura: los intercambios se
registran por hora, se liquidan al cierre del mes junto con la factura, y el
piso de cada vendedor es la permuta antes de su hora de corte y el precio de
bolsa de la hora a partir de ella. Como la hora de corte depende de la
importación total, solo se conoce al cerrar el mes; en una liquidación hecha
sobre datos medidos eso equivale a suponer conocidos los totales del mes, y
así se declara. En la frontera principal ningún mes alcanza la hora de corte;
en la secundaria, el reparto del exceso entre las horas apenas altera su
valor.

\subsection{El precio del excedente que supera el crédito}
\label{sub:preg-mcm}

Tampoco queda abierta. La regla definitiva de la Resolución CREG 101 072 de
2025 valora ese excedente a la variable MCm, el costo promedio ponderado de los
contratos con destino al mercado regulado. Pero el parágrafo del artículo 25
de la Resolución CREG 174, en la redacción que le dio la Resolución CREG 101 087
de 2025, mantiene el precio de bolsa hasta que la Comisión expida la
metodología de traslado, y los conceptos CREG 3018 y 3023 de abril de 2026
confirman que la transición sigue vigente. En todo el horizonte del trabajo el
excedente sobrante se valora, por tanto, al precio de bolsa de cada hora. Es lo
que hace la autogeneración individual del modelo; el colectivo usaba la media
mensual de la bolsa, que sobrevalora ese excedente porque las horas solares son
las de bolsa más barata, y debe corregirse.

\subsection{Con qué porcentaje reparte el colectivo}
\label{sub:preg-pde}

El artículo 19 de la Resolución CREG 101 072 deja el porcentaje de
distribución de los excedentes a lo que acuerden los integrantes. El modelo
lo calcula en proporción a la generación medida, que en la frontera
principal depende de lo que cada medidor alcanza a ver; con las placas
idénticas, repartir por capacidad instalada daría la misma proporción a las
cinco. Según la respuesta cambian las cuotas de cada institución y, con
ellas, a quién le conviene el colectivo.

\subsection{Si la autogeneración remota se modela como compensación}
\label{sub:preg-agr}

La Resolución CREG 101 099 de 2026 no compensa. La generación se entrega al
mercado mayorista a través de un agente generador, el consumo se atiende
como usuario no regulado con un contrato registrado ante el ASIC, y el
reparto horario con porcentaje fijo sirve solo para devolver el cargo por
confiabilidad. El escenario la modela como una compensación horaria entre
fronteras. Queda por decidir si se reescribe con esos artículos, se conserva
como referencia prospectiva o se retira, teniendo presente que la revisión
externa consideró el régimen cerrado a esta comunidad.

\placeholderfig{Faltan por trasladar aquí las decisiones abiertas que hoy
viven en el registro de correcciones.}
